\subsubsection*{Summary and Transition}

\begin{remark*}[Summary]
Functional completeness measures the expressive strength of a connective basis.
A functionally complete basis can represent every finite Boolean function by a
formula using only the connectives in that basis. Normal forms supply the main
construction: truth-table rows can be converted into formulas, and those
formulas can be assembled into a representative for the desired Boolean
function.
\end{remark*}

\begin{remark*}[Section map]
The standard connectives \(\neg\), \(\land\), and \(\lor\) are functionally
complete because they can build disjunctive normal forms. De Morgan laws reduce
this basis to \(\{\neg,\land\}\) or \(\{\neg,\lor\}\). NAND and NOR go further:
each single connective can define a complete two-connective basis and therefore
can represent every Boolean function.
\end{remark*}

\begin{remark*}[Transition]
The next section moves from expressive power to argument forms and informal
proof patterns. Functional completeness explains which truth functions can be
expressed. Argument forms ask which patterns of inference preserve truth from
premises to conclusions.
\end{remark*}
